\startsection[title={Sensores}]
	A princípio, as variáveis medidas podem apresentar uma colinearidade relativamente alta, mas o método de analise para produção de \getbuffer[sens] é projetada para funcionar através da utilização de uma caixa preta (LSTM), portanto o projeto do sistema de medição tem por papel medir todas as variáveis ao alcançe dos sensores disponíveis.

	\startitemize
		\starthead{Temperatura}
			A maior acurácia de medição, dentre os disponíveis, é proporcionado pelo sensor de temperatura. A \in{Figura}[fig:thermostate] ilustra o sensor utilizado, este opera com o protocolo {\emph one-wire}, ou seja, somente uma conexão é necessária para a comunicação com o controlador.

			\startplacefigure[title={Sensor de temperatura DS18B20}, reference={fig:thermostate}]
				\externalfigure[image/thermostate.png][height=400pt]
			\stopplacefigure

			A faixa de medição desse sensor é entre \math{\left\lbracket -55, 120 \right\rbracket \celsius}, sendo a precisão, dentro dessa faixa, de \math{\pm 0.5}.
		\stophead

		\starthead{Potência de Oxirredução}
			O potêncial de oxidação-redução da solução é uma medida representativa so equilibrio iônico desta. Esse sensor atua como um substituto ao sensor de pH.

			\startplacefigure[title={Sensor de potêncial de oxirredução}, reference={fig:oxireduction_potential}]
				\externalfigure[image/oxireduction_potential.png][height=400pt, width=\textwidth]
			\stopplacefigure

			A faixa de medição desse sensor é entre 
		\stophead

		\starthead{Hidrogênio}
			O papel do sensor de hidrogênio é de permitir o treinamento do algorítimo LSTM. Para a predição da produção de hidrogênio no processo fermentativo.
		\stophead

		\starthead{Metano}
			Outro sensor que serve como base para o treinamento do soft sensor para a substituição do sensor físico, neste caso, de metano.
		\stophead

		\starthead{Ácidos Graxos Voláteis}
			Os ácidos gráxos, são importantes para a produção de hidroênio, por conta da sua utilizasão em eletrólise microbiana, permitindo a existência de um processo secundário de evolução de hidrogênio. Portanto a sua determinação é importante para trabalhos futuros de otimização da produção de hidrogênio.
		\stophead

		%\starthead{Substrato}
		%	O substrato é formado por uma composição complexa e existem algumas dificuldades em medir sua concentração {\emph online}, portanto, a presença do substrato impoe a necessidade de uma amostragem recorrente, e medida {\emph offline}.
		%\stophead

		%\starthead{Biomassa}
		%	Assim como o substrato, a biomassa tem suas particularidades, composição complexa e tem de ser medida com cromatografia {\emph offline}.
		%\stophead
	\stopitemize
\stopsection
